\subsection{MATH 选中程序清单}
\label{app:panel-mechanisms}
表~\ref{tab:panel-mechanisms} 说明重复研究中选定的九个冻结成员。H1--H8 是生成成员的中性编号，生成种子均为 $0$。所列操作表示各程序实现了什么，不表示每个分支均被触发或提高了正确性。平均调用数的准入上限不是单次执行硬上限，一次逻辑调用也可能请求多个样本。

\begin{table}[htbp]
\centering\small
\begin{tabular}{@{}p{0.29\linewidth}p{0.66\linewidth}@{}}
\toprule
冻结成员 & 源代码定义的执行路径 \\
\midrule
\bare{} & 一次温度 $0$ 的求解，按规定格式给出最终答案。 \\
H1 (DeepSeek) & 三次温度 $0.4$ 的求解；归一化答案后投票；无答案或并列时采用贪心回退。 \\
H2 (ERNIE) & 起草解答，再让求解器验证或修正提取的答案；返回检查后答案。 \\
H3 (GLM) & 最多五种提示、温度 $0$ 的求解；三者一致即停止；分歧未解决时裁决，并设有格式回退。 \\
H4 (Kimi) & 一次调用请求五个温度 $0.7$ 的样本；提取答案后投票。 \\
H5 (Kimi) & 初次求解、验证与另一种求解；投票，没有多数时额外调用以解决分歧。 \\
H6 (MiniMax) & 草稿、批评；不一致时增加平局裁决调用；解析失败时使用此前可用答案。 \\
H7 (MiniMax) & 一次调用请求五个温度 $0.7$ 的样本；归一化并投票，并列时取最先出现者。 \\
H8 (Qwen) & 草稿加独立求解与审核；比较答案，冲突时请求裁决，并设有格式修复路径。 \\
\bottomrule
\end{tabular}
\caption{\textbf{重复研究具体选中了哪些程序。}所有调用使用同一冻结求解器；验证和裁决均为求解器提示，不能访问基准标签。六个席位按构建器分层选出，H5（Kimi）与 H7（MiniMax） 为两个不分层席位。}
\label{tab:panel-mechanisms}
\end{table}
